Cross correlation - 
Two signals are overlayed and moved across each other, one signal stays
static while the other signal is "moved" across the other one. The "correlation" 
at each sample from rest the moving signal has been moved is measured, known as 
the lag. There will be a peak in correlation at the lag at which both signals 
are the most similar.
This lag, given in samples, tells us how many samples two identical signals are
apart. This is a very powerful calculation which we can use in this context to
determine the time difference it takes for an identical signal to arrive at each
microphone. 
In a binaural context, the lag will be the Interaural Time Difference. The size 
and shape of the head will affect this, as well as the position of the source in 
relation to the head.
In a microphone array context, the lag between each pair of microphones can be
calculated, resulting in the Time Difference of Arrival.


Reverb Correction - 
While the initial measurements for this experiment were taken in a "dead" room,
a studio room in Maynooth University, in order for the experiment to be reproduceible,
correction for reverberation and reflections needs to be considered.
The easiest way of correcting for this is to window the recorded signal to only
include the direct sound, removing any early/late reflections.
Two geometric methods were investigated for this based on the geometry of the room
initial recordings were taken in.

The first method uses the room radius. The room radius within a room is the 
point at which the reverberant and direct sound fields are equal. Within the
room radius, direct sound dominates, beyond the room radius, the reverberant
field dominates.

The second method uses the dimensions of the room, and coordinates of the
source in order to determine the first possible re

These methods, while they work for this room and this set of measurements, would 
not be able to be generalised for all environments. A more appropriate method may
be to select a standard window size, or set of window sizes that the user can 
select from depending on their current environment, for example, highly 
reverberant rooms may require much shorter window lengths than less reverberant 
rooms, which may get away with the use of longer window lengths.

******** Need to compare window length differences in terms of DOA Estimation ****